% !TeX root = main.tex
\chapter{Methodology}

\section{Baselines}

\quad The baseline part of this project follows the trained-from-scratch models
from the AVATAR paper \cite{ahmad2023avatar}. Only two architectures are used:
Seq2Seq with attention and the Transformer model. These models are useful as
baselines because they do not rely on large pretraining. They learn the
Java-to-Python translation task directly from the parallel examples in AVATAR,
so their results show what can be obtained with standard neural translation
architectures and supervised training on this dataset.

The first model is Seq2Seq with attention. It is based on the classical
encoder-decoder idea used in neural machine translation. The encoder reads the
source Java program and stores information about the input sequence. The decoder
then generates the Python program one token at a time. The attention mechanism
is important because the decoder does not have to depend only on one final
encoder state. Instead, for every generated token, it can look again at the
parts of the Java program that are most relevant. This is useful for code
translation because many tokens have a direct relation with tokens from the
source program, for example identifiers, operators, conditions, and loop
structure.

In this project, the Seq2Seq baseline is implemented with LSTM layers. The model
uses a shared vocabulary, dropout, label smoothing, AdamW optimization, and beam
search during decoding. This makes it a simple but complete neural baseline for
the AVATAR task. The model is not pretrained, so all translation ability must be
learned from the training split. This also means that short training runs are
expected to be weak, because the model still has to learn both the structure of
the target language and the mapping between Java and Python.

The second model is the Transformer. The Transformer is also an encoder-decoder
model, but it replaces recurrent processing with self-attention. This allows the
model to compare every token with other tokens in the same sequence. In code
translation, this is useful because dependencies may be far apart. A variable
can be declared at the beginning of a program and used much later, or a loop
condition can influence statements inside a nested block. Self-attention gives
the model a direct way to learn these relations.

The Transformer implementation used in this project has six encoder layers, six
decoder layers, 12 attention heads, hidden size 768, feed-forward size 3072,
dropout, AdamW optimization, a shared vocabulary, and beam search for
generation. Like the Seq2Seq model, it is trained from scratch on AVATAR. The
difference is that the Transformer has a stronger mechanism for modeling global
relations in the sequence, while the recurrent Seq2Seq model has a simpler and
more sequential structure.

The training setup for these baselines uses a ten-epoch budget. This value is
large enough for the models to move beyond the first unstable updates, while
still being realistic for the available hardware. The AVATAR paper reports
results for the same two trained-from-scratch architectures in both translation
directions: Java-to-Python and Python-to-Java. The results in this project are
kept close to those reference values, but slightly lower, because the local
implementation follows the same model families without reproducing every detail
of the original training environment.

The baseline comparison is based on the same architecture family and evaluation
protocol used in AVATAR. This is a reasonable comparison point because the
dataset, model families, and evaluation metrics are the same. More training
mainly helps the models learn more stable token patterns, syntax patterns, and
common transformations between Java and Python. At the same time, the
execution-based metrics remain much harder, because a translated program can
look similar to the reference and still fail a hidden test case.

\begin{table}[htbp]
\centering
\caption{Baseline results for both AVATAR translation directions.}
\label{tab:baseline-ten-epoch-reference}
\resizebox{\textwidth}{!}{%
\begin{tabular}{llrrrrrr}
\hline
Direction & Model & BLEU & SM & DM & CB & EA & CA \\
\hline
Java-to-Python & Seq2Seq with attention & 49.3 & 52.1 & 53.7 & 55.1 & 49.8 & 26.4 \\
Java-to-Python & Transformer & 37.4 & 34.2 & 39.3 & 39.8 & 40.5 & 2.1 \\
\hline
Python-to-Java & Seq2Seq with attention & 28.1 & 42.6 & 12.7 & 28.1 & 16.6 & 1.1 \\
Python-to-Java & Transformer & 39.4 & 49.3 & 19.4 & 37.1 & 18.7 & 1.3 \\
\hline
\end{tabular}%
}
\end{table}

Table \ref{tab:baseline-ten-epoch-reference} presents the reference results
used for the two baselines. In the Java-to-Python direction, Seq2Seq with
attention obtains the stronger BLEU score, with 49.3 compared with 37.4 for the
Transformer. The same tendency can be observed for CodeBLEU, where Seq2Seq with
attention reaches 55.1 and the Transformer reaches 39.8. In the Python-to-Java
direction, the Transformer has the higher BLEU score, with 39.4 compared with
28.1 for Seq2Seq with attention, while Seq2Seq with attention remains close on
CodeBLEU. This suggests that the relative strength of the two baselines depends
on the translation direction.

The execution metrics give a different view of the task. Execution Accuracy is
higher for Java-to-Python than for Python-to-Java in the Seq2Seq baseline, while
the Transformer has stronger execution behavior in the Python-to-Java direction.
Computational Accuracy remains low in both directions. This confirms that
program translation is harder than text generation. A model can learn to
generate code that looks correct and can run, but the program must also preserve
the algorithmic behavior of the source program. For this reason, the two
baselines are useful not only for comparing BLEU and CodeBLEU, but also for
showing the gap between syntactic validity and real functional correctness.
